\documentclass[a4paper,twoside]{article}

\usepackage{morewrites}

\usepackage[svgnames,dvipsnames,x11names]{xcolor}
\usepackage{graphicx}
\usepackage{texsmith-lists}
\usepackage[english]{babel}
\usepackage[colorlinks=true, linkcolor=NavyBlue, urlcolor=NavyBlue, citecolor=NavyBlue]{hyperref}
\usepackage[a4paper]{geometry}
\usepackage{ts-fonts}
\usepackage{amsmath}
\usepackage{float}
\usepackage{seqsplit}
\usepackage{ts-code}

\usepackage{lastpage}
\usepackage{refcount}
\makeatletter
\def\ps@plain{
\let\@oddhead\@empty
\let\@evenhead\@empty
\def\@oddfoot{
\hfil
\ifnum\getpagerefnumber{LastPage}>1\relax
\thepage
\fi
\hfil}
\let\@evenfoot\@oddfoot
}
\makeatother

\title{Build}
\author{}\date{}
\begin{document}
\maketitle
To build your \LaTeX{} document into a PDF, use the \tscodeinline{-\allowbreak{}-\allowbreak{}build} (or \tscodeinline{-\allowbreak{}b}) option. This command compiles the generated \LaTeX{} file using the default engine. TeXSmith uses \tscodeinline{tectonic} by default as it is the only self-contained \LaTeX{} engine, but you can configure it to use \tscodeinline{lualatex} or \tscodeinline{xelatex} via \tscodeinline{latexmk} if preferred. Here is an example of a simple document built with different engines:
\begin{tscode}[lang=bash, engine=pygments]
\begin{Verbatim}[commandchars=\\\{\},breaklines, breakanywhere, commandchars=\\\{\}]
\PYZdl{}\PY{+w}{ }uv\PY{+w}{ }run\PY{+w}{ }texsmith\PY{+w}{ }cheese.md\PY{+w}{ }cheese.bib\PY{+w}{ }\PYZhy{}tarticle\PY{+w}{ }\PYZhy{}\PYZhy{}build\PY{+w}{ }\PYZhy{}\PYZhy{}engine\PY{o}{=}xelatex
\PYZdl{}\PY{+w}{ }uv\PY{+w}{ }run\PY{+w}{ }texsmith\PY{+w}{ }cheese.md\PY{+w}{ }cheese.bib\PY{+w}{ }\PYZhy{}tarticle\PY{+w}{ }\PYZhy{}\PYZhy{}build\PY{+w}{ }\PYZhy{}\PYZhy{}engine\PY{o}{=}lualatex
\PYZdl{}\PY{+w}{ }uv\PY{+w}{ }run\PY{+w}{ }texsmith\PY{+w}{ }cheese.md\PY{+w}{ }cheese.bib\PY{+w}{ }\PYZhy{}tarticle\PY{+w}{ }\PYZhy{}\PYZhy{}build\PY{+w}{ }\PYZhy{}\PYZhy{}engine\PY{o}{=}tectonic
\end{Verbatim}
\end{tscode}
\end{document}
